\section{USAJMO 2020}
        \begin{enumerate}
        	\item (\textit{USAJMO 2020}) \textbf{Note:} For any geometry problem whose statement begins with an asterisk $(*)$, the first page of the solution must be a large, in-scale, clearly labeled diagram. Failure to meet this requirement will result in an automatic 1-point deduction.


 Let $n \geq 2$ be an integer. Carl has $n$ books arranged on a bookshelf. Each book has a height and a width. No two books have the same height, and no two books have the same width. Initially, the books are arranged in increasing order of height from left to right. In a move, Carl picks any two adjacent books where the left book is wider and shorter than the right book, and swaps their locations. Carl does this repeatedly until no further moves are possible.  Prove that regardless of how Carl makes his moves, he must stop after a finite number of moves, and when he does stop, the books are sorted in increasing order of width from left to right.


 

	\item (\textit{USAJMO 2020}) \textbf{Note:} For any geometry problem whose statement begins with an asterisk $(*)$, the first page of the solution must be a large, in-scale, clearly labeled diagram. Failure to meet this requirement will result in an automatic 1-point deduction.


 Let $\omega$ be the incircle of a fixed equilateral triangle $ABC$. Let $\ell$ be a variable line that is tangent to $\omega$ and meets the interior of segments $BC$ and $CA$ at points $P$ and $Q$, respectively. A point $R$ is chosen such that $PR = PA$ and $QR = QB$. Find all possible locations of the point $R$, over all choices of $\ell$.


 

	\item (\textit{USAJMO 2020}) \textbf{Note:} For any geometry problem whose statement begins with an asterisk $(*)$, the first page of the solution must be a large, in-scale, clearly labeled diagram. Failure to meet this requirement will result in an automatic 1-point deduction.


 An empty $2020 \times 2020 \times 2020$ cube is given, and a $2020 \times 2020$ grid of square unit cells is drawn  on each of its six faces. A beam is a $1 \times 1 \times 2020$ rectangular prism. Several beams are placed inside the cube subject to the following conditions:


 
                \begin{itemize}
                    \item The two $1 \times 1$ faces of each beam coincide with unit cells lying on opposite faces of the cube. (Hence, there are $3 \cdot {2020}^2$ possible positions for a beam.)
\item No two beams have intersecting interiors.
\item The interiors of each of the four $1 \times 2020$ faces of each beam touch either a face of the cube or the interior of the face of another beam.
                \end{itemize}


                What is the smallest positive number of beams that can be placed to satisfy these conditions?


 

	\item (\textit{USAJMO 2020}) \textbf{Note:} For any geometry problem whose statement begins with an asterisk $(*)$, the first page of the solution must be a large, in-scale, clearly labeled diagram. Failure to meet this requirement will result in an automatic 1-point deduction.


 Let $ABCD$ be a convex quadrilateral inscribed in a circle and satisfying $DA < AB = BC < CD$. Points $E$ and $F$ are chosen on sides $CD$ and $AB$ such that $BE \perp AC$ and $EF \parallel BC$. Prove that $FB = FD$.


 

	\item (\textit{USAJMO 2020}) \textbf{Note:} For any geometry problem whose statement begins with an asterisk $(*)$, the first page of the solution must be a large, in-scale, clearly labeled diagram. Failure to meet this requirement will result in an automatic 1-point deduction.


 Suppose that $(a_1,b_1),$ $(a_2,b_2),$ $\dots,$ $(a_{100},b_{100})$ are distinct ordered pairs of nonnegative integers. Let $N$ denote the number of pairs of integers $(i,j)$ satisfying $1\leq i<j\leq 100$ and $|a_ib_j-a_jb_i|=1$. Determine the largest possible value of $N$ over all possible choices of the $100$ ordered pairs.


 

	\item (\textit{USAJMO 2020}) \textbf{Note:} For any geometry problem whose statement begins with an asterisk $(*)$, the first page of the solution must be a large, in-scale, clearly labeled diagram. Failure to meet this requirement will result in an automatic 1-point deduction.


 Let $n \geq 2$ be an integer. Let $P(x_1, x_2, \ldots, x_n)$ be a nonconstant $n$-variable polynomial with real coefficients. Assume that whenever $r_1, r_2, \ldots , r_n$ are real numbers, at least two of which are equal, we have $P(r_1, r_2, \ldots , r_n) = 0$. Prove that $P(x_1, x_2, \ldots, x_n)$ cannot be written as the sum of fewer than $n!$ monomials. (A monomial is a polynomial of the form $cx^{d_1}_1 x^{d_2}_2\ldots x^{d_n}_n$, where $c$ is a nonzero real number and $d_1$, $d_2$, $\ldots$, $d_n$ are nonnegative integers.)


 

\end{enumerate}